\chapter{Checking results and copies}
\label{ch:checking}

Chapter~\ref{ch:classification} sorts results into kinds, and
Chapter~\ref{ch:provenance} describes the records behind them.  This
chapter gives the commands that check them:
\begin{itemize}
\item the exam, which reruns every proof of the library
  (Section~\ref{sec:vnb-test});
\item the REPL commands that show the classification of one result
  (Section~\ref{sec:status-query});
\item the commands for signed trees (Section~\ref{sec:signed-trees-use})
  and for signed releases (Section~\ref{sec:signed-releases-use}).  What
  the two signatures are and what they protect is in
  Sections~\ref{sec:signed-trees} and~\ref{sec:signed-releases}.
\end{itemize}

%% =========================================================

\section{The exam}
\label{sec:vnb-test}

The \emph{exam} is the standing examination of the \textbf{theorem}
bucket.  It is a full library build in which every proof script runs,
every proof page is typed back in and must reach a grounded graph, and
every certificate in \texttt{certificates/} is rewritten from the run.
Definitions and PSS members are excused by construction: the former have
nothing to prove, the latter are accepted on a warrant.

\begin{VNBsnip}
vnb-band WORKER --build --exam
\end{VNBsnip}

\noindent
runs it on a worker machine; \texttt{vnb-nightly-exam} runs it every
night, skips a tree that the last exam already passed, and records the
outcome in \texttt{last-exam.json}.  The everyday build is the
\emph{certified} build: a theorem whose certificate is still valid is
installed without running its proof, and only the theorems whose
statement, citations or definitions changed since the exam are proved
again.  A change to a kernel file invalidates every certificate, so the
everyday build after such a change is itself an exam.

\begin{remark}
The exam is \emph{not} the regression suite.  \texttt{test-suite.scm}
checks that the machinery behaves; the exam checks that the library's
theorems still derive.  The two answer different questions and both are
run.
\end{remark}

\begin{remark}
The script \texttt{./vnb-test} and its stamp
\nolinkurl{reference/certification.scm} are not maintained; the exam is
run by the command above.
\end{remark}

%% =========================================================
\section{Inspecting a result}
\label{sec:status-query}

Three commands expose the classification at the REPL.

\begin{itemize}
\item \texttt{(status 'name)} prints, on one screen, a result's
  provenance, whether it is a PSS member, any warrant it carries, and the
  bill.

\item \texttt{(status-audit)} writes \texttt{reference/STATUS-AUDIT.md}: a
  census of the whole library by kind, flagging hygiene issues --- asserted
  results with no warrant, theorem-shaped names that have drifted out of
  the theorem bucket, and the like.

\item \texttt{(pss-letter 'name)} drafts a plain-English ``recommendation
  letter'' for a candidate PSS member: an English rendering of what it
  says, its current standing, and an assessment against the admission
  gates (true; a sound, terminating rewrite shape; general rather than
  bespoke; non-redundant).  It is an aid to deciding whether a result
  belongs in the curated support set, not an automatic admission.
\end{itemize}


%% =========================================================
\section{Using signed trees}
\label{sec:signed-trees-use}

This section gives the commands for the signed trees of
Section~\ref{sec:signed-trees}.

\paragraph{The pieces.}  Two scripts, \texttt{vnb-trust} and
\texttt{vnb-start}, whose master copies are in the tree under
\texttt{trust/}.  The user copies them once, after checking the
tree, into \texttt{\textasciitilde/bin}, outside the tree, so that a
later change to the tree cannot change them too:
\begin{verbatim}
  cp trust/vnb-trust trust/vnb-start ~/bin/
\end{verbatim}
\begin{itemize}
\item \texttt{vnb-trust} has the commands \texttt{setup KEY} (once),
  \texttt{accept DIR}, \texttt{verify DIR}, \texttt{list} and
  \texttt{release-sign} (the last is for an author signing a
  release).
\item \texttt{vnb-start} runs \texttt{vnb-trust verify} on the tree and
  starts VNB only if it succeeds.  On a mismatch it lists the changed
  files and offers to stop, to fetch a fresh copy, or to start anyway
  once the word \texttt{unverified} is typed in full.
\end{itemize}
The signature is an SSH signature (\texttt{ssh-keygen -Y sign}) in its
own namespace, \texttt{vnb-tree}, so it cannot be mistaken for any
other signature made with the same key.  The hashes are SHA-256.

\paragraph{The manifest.}  To \emph{accept} a tree is to run
\texttt{vnb-trust accept DIR} on it.  The command hashes the files,
shows a summary and asks \texttt{Accept and sign this tree? [y/N]}.  On
\texttt{y}, and the key's passphrase, it writes a \emph{manifest} to
\texttt{\textasciitilde/vnb-signed/}, with its signature beside it: a
few header lines, then one line per file, its SHA-256 hash and its
path.  Every file of the tree is
listed except
\begin{itemize}
\item \texttt{.git/}, \texttt{scratch/}, \texttt{scratchpad/}, and the
  files \texttt{*.bci}, \texttt{*.ext}, \texttt{*.elc};
\item the documents under \texttt{docs/} and \texttt{reference/}, which
  VNB never loads.  The \texttt{.sexp} and \texttt{.scm} files there are
  listed, because the load reads some of them.
\end{itemize}
So the source, the compiled files, the band, the certificates and the
Emacs files are all covered.  A session of VNB does not change any
listed file, so an accepted tree stays accepted while it is used.  The
manifest and its signature stay on the user's computer.

\paragraph{Updating.}  A new copy, a new release or an update fetched
from elsewhere, is compared in the same way and, until the user
accepts it, is refused with the list of files that differ.  The user
then checks it by one of the three ways of Section~\ref{sec:signed-trees}
and accepts it, or keeps
the old copy.

Two rules keep these signatures meaningful.  The copies of \texttt{vnb-trust} and
\texttt{vnb-start} in \texttt{\textasciitilde/bin} must not be replaced
without reading the change.  And the signing key should have a
passphrase.  The setup takes a few commands; they are in the
repository's \texttt{README.md}, under ``Keeping your copy checked''.

\paragraph{For work that must withstand an adversary.}  The steps are
those above, with every choice made the strict way.  Three things
differ from ordinary use.
\begin{itemize}
\item The copy is a copy of a published release, and its signature is
  checked (Section~\ref{sec:signed-releases}) with the author's key
  compared in both places it is published, on GitHub and in the
  author's Bluesky profile.  Ordinary use may
  accept a copy on lighter grounds.
\item Before accepting, the user runs the exam on the user's own
  computer:
\begin{verbatim}
  VNB_CERTIFIED=off ./prover --build-band
\end{verbatim}
  It reruns every proof through the kernel and rebuilds the
  certificates and the band.  This is the step that matters most.  In
  ordinary use the certificates and the band come from the author;
  after the exam, the user relies on nothing the author or the
  author's server did, only on the kernel, whose source is in the
  tree.
\item VNB is started only with \texttt{vnb-start}, never directly.
\end{itemize}
The limits listed under ``What it protects, and what it does not''
(Section~\ref{sec:signed-trees}) still apply.


%% =========================================================
\section{Making and checking a release}
\label{sec:signed-releases-use}

This section gives the commands for the signed releases of
Section~\ref{sec:signed-releases}.

\paragraph{How the author makes one.}  The order is:
\begin{enumerate}
\item the author checks the tree, by the exam, and accepts it
  (Section~\ref{sec:signed-trees-use});
\item on the machine that holds the repository, \texttt{vnb-release-manifest} \textit{VERSION}
  writes the manifest straight from git's object store;
\item on the author's computer, \texttt{vnb-trust release-sign}
  compares every file of the manifest with the same file in the
  accepted tree, and signs only if all agree.  Files under
  \texttt{scratch/} and \texttt{scratchpad/}, which VNB does not load,
  are counted and shown, not compared;
\item the signature is copied back, both files are committed, and the
  release commit is tagged and pushed.
\end{enumerate}
So a release is exactly a tree the author accepted.

\paragraph{How anyone checks one.}  With standard tools only:
\texttt{ssh-keygen -Y verify} for the signature, then
\texttt{git ls-files -z | xargs -0 sha256sum} at the release commit,
compared with the manifest.  The commands are in the repository's
\texttt{README.md}, under ``Verifying a release''.

